\anexo{Prompts Utilizados na Estratégia Map-Reduce}
\label{an:prompts-map-reduce}

Este anexo apresenta os \textit{prompts} utilizados na abordagem baseada no paradigma \textit{Map-Reduce} para a extração estruturada de informações a partir das transcrições das reuniões.  
Nessa estratégia, dois tipos de \textit{prompts} são empregados:  
(i) o \textbf{prompt de mapeamento}, utilizado para processar individualmente segmentos da transcrição, e  
(ii) o \textbf{prompt de redução}, responsável por consolidar as saídas parciais em uma síntese unificada que servirá como base para a construção final da ata.

\section*{Prompt de Mapeamento (Map)}
\addcontentsline{toc}{subsection}{Prompt de Mapeamento (Map)}

\begin{lstlisting}[language=Python]
Você é um agente de IA responsável por analisar segmentos de uma transcrição
de reunião do Tribunal Regional Eleitoral do Ceará.
Sua tarefa é extrair, APENAS a partir do segmento fornecido, as seguintes
informações:

1. Participantes e seus cargos
   - Identifique todos os participantes mencionados.
   - Sempre que possível, associe cada participante ao seu cargo ou função.

2. Pautas/Assuntos discutidos
   - Identifique as pautas tratadas no segmento.
   - Quando possível, associe a pauta ao participante ou órgão que a propôs.
   - Formate como: [Título] – [Descrição].

3. Deliberações tomadas
   - Identifique decisões, encaminhamentos ou ações aprovadas.
   - Forneça o tema associado e descreva a ação tomada.

4. Discussões e contexto
   - Identifique os temas debatidos.
   - Descreva brevemente o contexto e os pontos essenciais discutidos.

IMPORTANTE:
- A saída deve conter listas claras e separadas para cada um dos quatro itens.
- Não adicione introduções, resumos ou interpretações externas.
- Mantenha o texto inteiramente em Português brasileiro.

SEGMENTO DA TRANSCRIÇÃO:
{chunk_text}
\end{lstlisting}


\section*{Prompt de Redução (Reduce)}
\addcontentsline{toc}{subsection}{Prompt de Redução (Reduce)}

\begin{lstlisting}[language=Python]
Você é um agente de consolidação responsável por unificar e limpar informações
provenientes de múltiplos segmentos de transcrição.

Sua tarefa é produzir uma síntese final organizada por seções estruturais,
como:
- Participantes
- Pautas
- Deliberações
- Temas discutidos

Instruções:
- Combine e consolide todos os dados fornecidos.
- Elimine duplicações, ruídos e informações contraditórias.
- Mantenha apenas o que estiver presente nos segmentos.
- A saída deve ser um texto único, coeso e completamente estruturado.
- O resultado deve estar em Português brasileiro.

DADOS BRUTOS CONCATENADOS:
{dados_brutos}
\end{lstlisting}